\documentclass{article}
\usepackage{amsmath,amssymb}
\usepackage{xurl}
\usepackage[hidelinks]{hyperref}
% Shared commands for notes in docs/paper-gaps/.

\DeclareUnicodeCharacter{2080}{\ifmmode{}_{0}\else\textsubscript{0}\fi}
\DeclareUnicodeCharacter{2081}{\ifmmode{}_{1}\else\textsubscript{1}\fi}
\DeclareUnicodeCharacter{2082}{\ifmmode{}_{2}\else\textsubscript{2}\fi}
\DeclareUnicodeCharacter{2099}{\ifmmode{}_{n}\else\textsubscript{n}\fi}

\newcommand{\Lean}{\textsc{Lean}}
\ExplSyntaxOn
\NewDocumentCommand{\leanid}{m}{
  \ifmmode
    \textsf{\detokenize{#1}}
  \else
    \group_begin:
    \urlstyle{sf}
    \tl_set:Nn \l_tmpa_tl {#1}
    \tl_replace_all:Nnn \l_tmpa_tl {\_} {_}
    \exp_args:NV \path \l_tmpa_tl
    \group_end:
  \fi
}
\ExplSyntaxOff
\providecommand{\tr}{\operatorname{tr}}
\newcommand{\tnleanIssueBase}{https://github.com/LionSR/TNLean/issues/}
\newcommand{\tnleanPrBase}{https://github.com/LionSR/TNLean/pull/}
\newcommand{\tnleanDocBase}{https://sirui-lu.com/TNLean/docs/}
\newcommand{\ghissue}[1]{\href{\tnleanIssueBase#1}{GitHub issue \##1}}
\newcommand{\ghpr}[1]{\href{\tnleanPrBase#1}{PR \##1}}
\newcommand{\doclink}[1]{\href{\tnleanDocBase#1.html}{\path{#1}}}

% Dirac notation and the identity operator, as in blueprint/src/macros/common.tex.
% \providecommand keeps note-local definitions authoritative.
\usepackage{dsfont}
\providecommand{\ket}[1]{|#1\rangle}
\providecommand{\bra}[1]{\langle #1|}
\providecommand{\braket}[2]{\langle #1 | #2 \rangle}
\providecommand{\ketbra}[2]{|#1\rangle\!\langle #2|}
\providecommand{\Id}{\mathds{1}}
\providecommand{\one}{\mathds{1}}
\providecommand{\C}{\mathbb{C}}
\providecommand{\MN}[1]{M_{#1}(\C)}
\providecommand{\MD}{M_D(\C)}

% Verdict marker: \gapnote{<kind>}{<status>}. Kind is the policy
% classification (clarification, local-correction, scope-restriction,
% unfaithful, false-source, open-gap); status is open, wip (elimination
% underway), resolved, or historical. Machine-read by the site index;
% prints a verdict line.
\newcommand{\gapnote}[2]{\par\noindent\textbf{Verdict:} #1 (#2).\par}

\title{Physical bond separation and the remaining regular PEPS blocking step}
\author{}
\date{2026-10-05}
\begin{document}
\maketitle
\gapnote{scope-restriction}{open}

\section{Source and proved network statement}
The source is Schuch, Cirac, and P\'erez-Garc\'ia,
\href{https://arxiv.org/abs/1001.3807}{arXiv:1001.3807},
Observations 6.5--6.6, local source lines 1818--1915. The three figures
\path{renorm-gg-sym.pdf}, \path{renorm-g-id-and-split.pdf}, and
\path{renorm-fixedpoint.pdf} distinguish the adjacent bond separation
from the full $2\times2$ renormalization step.

Fix any finite group $G$, finite surplus index set $K$, and finite simple
graph $\Gamma$. Put $X=G\times G^K$. The local canonical averaging tensor
has one physical $X$ coordinate on every incident half-edge and one virtual
$X$ coordinate on every edge. Its group action translates all coordinates
simultaneously. The actual state sums the product of local tensors over
one virtual $X$ label per edge, with the group average divided by $|G|$
at every vertex.

The physical basis change $(a,b)\mapsto(a,a^{-1}b)$ is applied separately
to every half-edge at each vertex. If $D=|G^K|$, its exact output is
\begin{align}
 U\Psi_K=\Psi_0\otimes B_K
       =(\sqrt D)^{|E(\Gamma)|}\Psi_0\otimes\Omega_K,
 \qquad \|\Omega_K\|^2=1.
\end{align}
Here $\Psi_0$ is the actual single-regular-bond averaging network, and
$B_K$ is the product of endpoint equality vectors on all surplus bonds.
Every physical index and the full scalar are retained. In particular,
$K$ a singleton is the two-parallel-bond step in the source.

The proof expands the genuine tensor contraction, uses the explicit
regular relative-coordinate identity at each physical endpoint, and
extracts the equality factors from the group-label sum. It does not
assume the resulting global factorization. The map is a full-domain
physical isometry, and all physical overlaps are preserved.

\section{Exact restrictions and remaining work}
The canonical theorem has now been transported to arbitrary given physical
$G$-isometric site tensors on the bundled graph. The physical spaces and site
tensors may vary with the vertex. Local $G$-isometry gives positive factors
$c_v$ and the explicit normalized adjoints $J_v=c_v^{-1/2}T_v^*$.
They are isometric on the actual local physical ranges. Their product,
followed by the relative-coordinate unitary, restricts to a Hilbert-space
isometry $I$ on the product physical range, and the actual state is proved
to belong to that range. The exact state identity is
\begin{align}
 I\Psi_a=\left(\prod_v\sqrt{c_v}\right)
    (\sqrt D)^{|E(\Gamma)|}\Psi_0\otimes\Omega_K.
\end{align}
No surjectivity onto the original ambient physical spaces, supplied state
support, or supplied global equality is needed. Thus the physical-support
composition of Observation 6.4 is proved, with every local normalization
factor retained.

There is one uniform finite surplus set $K$ on every graph edge. The group
may be nonabelian. The graph need not be connected and may have isolated
vertices or be empty. Bundles are encoded by the product edge alphabet $X$;
no theorem about an arbitrary multigraph presentation or independently
varying edge bundle sizes is asserted. Inserted closure operators are not
part of this theorem.

The additional geometric step for Observation 6.6 is now derived on
untwisted doubled tori with coarse periods at least three. The actual
$2\times2$ contraction, including seams, gives the bundled coarse network;
its local G-isometry follows from the four fine-site hypotheses. The
surviving tensor is identified with the same original coarse tensor, and
both states are proved nonzero before normalization. Every positive scalar
is retained. See
\path{scp10_two_by_two_regular_fixed_point.tex} and
\leanid{TNLean.PEPS.exists_regularTwoByTwoOriginalTensorIsometry}.
The present bond-separation theorem alone still does not imply that
composition. Tiny-period physical Bell separation, inserted closures,
and other boundary conditions remain outside its scope, tracked in
\href{https://github.com/LionSR/TNLean/issues/8676}{issue 8676}.

\section{Declarations}
The production module is
\path{TNLean/PEPS/RegularGraphBondBlocking.lean}.
Its exact unnormalized and normalized contraction identities are
\leanid{TNLean.PEPS.regularGraphBlocking_averagingSite} and
\leanid{TNLean.PEPS.regularGraphBlocking_averagingSite_normalized}.
The physical isometry and its coefficient identity are
\leanid{TNLean.PEPS.regularGraphBlockingIsometry} and
\leanid{TNLean.PEPS.regularGraphBlockingIsometry_apply}.

The arbitrary-site extension is
\path{TNLean/PEPS/RegularGraphPhysicalBlocking.lean}. Its capstones are
\leanid{TNLean.PEPS.exists_regularGraphPhysicalDisentangler} and
\leanid{TNLean.PEPS.exists_regularGraphPhysicalSupportIsometry}.
The latter evaluates the actual Hilbert-space isometry on the original
state, using derived physical-range membership.
\end{document}
